\subsection{Principal coverings}

DEFINITION 2. --- Let B be a topological space and let G be a group. Endow G with the discrete topology. A principal fiber bundle with base B and group G is called a principal covering of B with group G.

This terminology is legitimate because such a B-space is a covering of B.

A principal G-bundle over a nonempty topological space has a degree, equal to Card G.

PROPOSITION 4. — Let B be a topological space, let (E, p) be a B-space, and let G be a discrete topological group acting on the right on E. The following assertions are equivalent:

\begin{enumerate}
\def\labelenumi{(\roman{enumi})}
\item
  The B-space E is a principal covering with group G;
\item
  For every \(g \in G\) and every \(x \in E\) , we have \(p(x \cdot g) = p(x)\) , the map p induces a homeomorphism from E/G onto B, and the map \(\theta: (x, g) \mapsto (x, x \cdot g)\) is a homeomorphism from E × G onto \(E \times_{B} E\);
\item
  The group G acts continuously and freely on E, the map p is étale, and its fibers are the orbits of G.
\end{enumerate}

(i)⇒(ii): This follows from I, p. 91 and from proposition 3 of I, p. 94. (ii)⇒(iii): Under hypothesis (ii), the action law of G is continuous, and the map p is surjective and open (TG, III, p. 10, lemma 2). Let \(x \in E\) ; the map \(\theta\) induces a bijection from \(\{x\} \times G\) onto \(\{x\} \times p^{-1}(p(x))\) , so the group G acts freely and its orbits are the fibers of p. If e denotes the identity element of G, the diagonal of \(E \times_{B} E\) is the image of \(E \times \{e\}\) by the homeomorphism \(\theta\). Since the group G is discrete, the set \(\{e\}\) is open in G, hence the diagonal \(\Delta_{E}\) is open in \(E \times_{B} E\). By Proposition 7 of I, p. 31, the map p is étale.

(iii)⇒(i): Since every fiber of p is an orbit of the action of G on E, the map p is surjective. As p is étale, for every point b of B, there exists an open neighborhood U of b and a continuous section s of p over U. The set \(s(\mathrm{U})\) is open in E and s induces a homeomorphism from U onto \(s(\mathrm{U})\) (I, p. 30, cor. 3). For every \(g \in G\), the set \(s(\mathrm{U}) \cdot g\) is open in E and the union of these sets over all \(g \in G\) is equal to \(p^{-1}(\mathrm{U})\). If g and \(g'\) are two distinct elements of G, the sets \(s(\mathrm{U}) \cdot g\) and \(s(\mathrm{U}) \cdot g'\) are disjoint, since G acts freely on E. The map \(f: U \times G \to p^{-1}(\mathrm{U})\) defined by \(f(u, g) = s(u) \cdot g\) is continuous for \((u, g) \in U \times G\) and is therefore a homeomorphism from \(U \times G\) onto \(p^{-1}(\mathrm{U})\), compatible with the right action of G. By definition, E is therefore a principal covering of B with group G.

Examples. --- 1) Let E be a topological space, G a discrete topological group acting continuously and freely on E on the right. For the canonical map \(p \colon \mathrm{E} \to \mathrm{E} / \mathrm{G}\) to make E a principal covering of E/G, it is necessary and sufficient that it be étale (condition (iii) of Proposition 4). For example, suppose there exists a topological space X and an étale map \(q \colon \mathrm{E} \to \mathrm{X}\) compatible with the action of G on E; then E is a principal covering of E/G. Indeed, let us denote by \(q' \colon \mathrm{E} / \mathrm{G} \to \mathrm{X}\) the map induced by \(q\); we have \(q = q' \circ p\), and therefore \(p\) is étale by Proposition 6, d) of I, p. 29.

\begin{enumerate}
\def\labelenumi{\arabic{enumi})}
\setcounter{enumi}{1}
\tightlist
\item
  Let \(q \colon \mathrm{E} \to \mathrm{X}\) be a separated étale morphism. The group \(\operatorname{Aut}_{\mathrm{X}}(\mathrm{E})\) , endowed with the discrete topology, acts continuously on the left on E. If the space E is connected, this action is free (I, p. 34, corollary 2). Let G denote the group \(\operatorname{Aut}_{\mathrm{X}}(\mathrm{E})^{\circ}\) and endow E with the right action opposite to that of \(\operatorname{Aut}_{\mathrm{X}}(\mathrm{E})\) . By the preceding example, the space E, endowed with the canonical map \(p \colon \mathrm{E} \to \mathrm{E} / \mathrm{G}\) , is a principal covering with group G.
\end{enumerate}

